\begin{theorem}[Integral--finite bridge]\label{mod:delta-integralfinitebridge}
Let $F$ be a nonarchimedean local field, let
$\chi:F^\times\to\mathbf C^\times$ have exact multiplicative conductor
$m\in\mathbf N$, let $\psi:F\to\mathbf C^\times$ have exact additive
conductor $n\in\mathbf Z$,
and let $\gamma\in F^\times$ be an admissible denominator, so that
$\operatorname{ord}_F(\gamma)=m+n$.  Let $du$ be an arbitrary positive
left Haar measure on the multiplicative group $U_F^0$.

Put
\[
 Q_m=U_F^0/U_F^m
\]
using exactly the finite quotient from the computational definition.  For
$z\in Q_m$, Lean defines $C_z$ without choosing a representative, as the
literal fibre over $z$ of the canonical quotient map $U_F^0\to Q_m$.  It
proves
\[
 U_F^0=\bigcup_{z\in Q_m}C_z,
 \qquad C_z\cap C_w=\varnothing\quad(z\ne w),
\]
and identifies $C_z$ with the left coset
$(\operatorname{out}z)U_F^m$ to prove the cover, its topology, and Haar
measure.  Every $C_z$ is open and Borel measurable.

For
\[
 g_\gamma(u)=\psi(u/\gamma)\chi(u)^{-1},
\]
membership $u\in C_z$ gives the exact identity
\[
 g_\gamma(u)=g^{\mathrm{fin}}_\gamma(z),
\]
where the right side is the existing representative-free
\texttt{finiteGaussSummand}.  Thus the additive argument has no inserted
minus sign, the quasi-character is inverted exactly once, and the same
admissible $\gamma$ is used on both sides.  The integrand is already proved
continuous, Borel measurable, and integrable on $U_F^0$; Lean also proves
its integrability on each measurable coset.

Left Haar invariance gives every coset the common measure
\[
 du(C_z)=du(U_F^m).
\]
The subgroup $U_F^m$ is open and contains $1$, while $U_F^0$ is compact.
Consequently Lean proves
\[
 0<du(U_F^m)<\infty,
 \qquad du(U_F^m)\ne0.
\]
Set
\[
 c=\bigl(du(U_F^m)\bigr)_{\mathbf R}.
\]
The conversion is justified by finiteness, and $c>0$, $c\ne0$, and
$\operatorname{ph}(c)=1$ are all proved explicitly.

The pointwise finite-coset decomposition is expressed as a finite sum of
indicator functions.  Each indicator set is measurable and each constant
term is proved integrable using the finite coset measure.  Termwise Bochner
integration then gives the literal identity
\[
 \int_{U_F^0}\psi(u/\gamma)\chi(u)^{-1}\,du
   =c\sum_{z\in Q_m}g^{\mathrm{fin}}_\gamma(z)
   =c\,\operatorname{finiteGaussSum}(\chi,\psi;\gamma).
\]
There is no factor involving $|Q_m|$, $q_F$, a sign, or a Haar
normalization.  The raw integral is nonzero because $c\ne0$ and the finite
Gauss sum was proved nonzero independently in
Definition~\ref{mod:delta-finitedefinition}; no local integral
nonvanishing theorem or First Main Lemma case is assumed.

Finally, positivity gives
$\operatorname{ph}(cG)=\operatorname{ph}(G)$.  Multiplying both sides by
the unchanged outside factor $\chi(\gamma)$ proves exactly
\[
 \Delta_F^{\mathrm{int}}(\chi,\psi;\gamma)
   =\Delta_F^{\mathrm{fin}}(\chi,\psi;\gamma).
\]

\LeanFile{LanglandsFirstMainLemma/Delta/IntegralFiniteBridge.lean}
\lean{LanglandsFirstMainLemma.integralFiniteCoset}
\lean{LanglandsFirstMainLemma.mem_integralFiniteCoset_iff}
\lean{LanglandsFirstMainLemma.integralFiniteCoset_eq_out_smul}
\lean{LanglandsFirstMainLemma.integralFiniteCoset_cover}
\lean{LanglandsFirstMainLemma.integralFiniteCoset_pairwiseDisjoint}
\lean{LanglandsFirstMainLemma.integralFiniteCoset_isOpen}
\lean{LanglandsFirstMainLemma.integralFiniteCoset_measurableSet}
\lean{LanglandsFirstMainLemma.unitHaarIntegrand_eq_on_integralFiniteCoset}
\lean{LanglandsFirstMainLemma.unitHaarIntegrand_integrableOn_integralFiniteCoset}
\lean{LanglandsFirstMainLemma.integralFiniteCoset_measure_eq}
\lean{LanglandsFirstMainLemma.integralFiniteCoset_measure_pos}
\lean{LanglandsFirstMainLemma.integralFiniteCoset_measure_lt_top}
\lean{LanglandsFirstMainLemma.integralFiniteCoset_measure_ne_zero}
\lean{LanglandsFirstMainLemma.integralFinitePositiveScalar}
\lean{LanglandsFirstMainLemma.integralFinitePositiveScalar_pos}
\lean{LanglandsFirstMainLemma.integralFinitePositiveScalar_ne_zero}
\lean{LanglandsFirstMainLemma.phase_integralFinitePositiveScalar}
\lean{LanglandsFirstMainLemma.unitHaarIntegrand_eq_coset_sum}
\lean{LanglandsFirstMainLemma.unitGaussIntegral_eq_positiveScalar_mul_finiteGaussSum}
\lean{LanglandsFirstMainLemma.unitGaussIntegral_ne_zero_of_finiteGaussSum}
\lean{LanglandsFirstMainLemma.deltaIntegral_eq_deltaFinite}
\leanok
\SourceText{Definition of Delta; Lemma lem:gamma-independent; Corollary cor:local-gauss-integral-nonzero.}
\uses{mod:delta-finitedefinition,mod:delta-integraldefinition}
\FormalizationContract The comparison is performed on the multiplicative
Haar space $U_F^0$ and uses exactly $U_F^0/U_F^m$.  Cosets are quotient-map
fibres, and quotient representatives occur only in a proved left-coset
description used for the cover, openness, and Haar invariance.  The common scalar is the
proved finite positive Haar measure of $U_F^m$.  The finite sum is the
existing unnormalised \texttt{finiteGaussSum}; the proof introduces no sign,
cardinality factor, normalization, nonvanishing postulate, bridge assumption,
or First Main Lemma case theorem.
\end{theorem}
